\subsection{Finite Observational Projection and Emergent Memory}

A central premise of the SBC framework is that the observable universe does not necessarily coincide with the total causal structure of physical reality.

Let \(\Psi_{\rm tot}(t)\) denote the state of the full causal structure, and let the observable cosmological state be defined as a projection,
\begin{equation}
\Psi_{\rm obs}(t)
=
\mathcal{P}\Psi_{\rm tot}(t),
\end{equation}
where \(\mathcal{P}\) is an observational projection operator associated with a finite observational window.

The complementary, unobserved sector is
\begin{equation}
\Psi_{\rm hid}(t)
=
(1-\mathcal{P})\Psi_{\rm tot}(t).
\end{equation}

If the total state evolves locally in the full causal space,
\begin{equation}
\frac{d\Psi_{\rm tot}}{dt}
=
\mathcal{L}\Psi_{\rm tot},
\end{equation}
then projection onto the observable sector yields an effective equation of the form
\begin{equation}
\frac{d\Psi_{\rm obs}}{dt}
=
\mathcal{P}\mathcal{L}\mathcal{P}\Psi_{\rm tot}
+
\mathcal{P}\mathcal{L}(1-\mathcal{P})\Psi_{\rm tot}.
\end{equation}

The second term represents the influence of the hidden causal sector on the observable dynamics. Since \(\Psi_{\rm hid}\) is not directly resolved within the observational window, its effect appears as a nonlocal memory contribution.

Formally, this can be written as
\begin{equation}
\frac{d\Psi_{\rm obs}}{dt}
=
\mathcal{L}_{\rm eff}\Psi_{\rm obs}(t)
+
\int_{-\infty}^{t}
K(t,t')\Psi_{\rm obs}(t')\,dt',
\end{equation}
where \(K(t,t')\) is an effective memory kernel induced by the eliminated hidden sector.

Thus, in the SBC framework, memory is not introduced as an independent assumption. It emerges from the finite observational projection:
\begin{equation}
{\rm finite\ observational\ window}
\quad
\Longrightarrow
\quad
{\rm incomplete\ state\ description}
\quad
\Longrightarrow
\quad
{\rm effective\ memory}.
\end{equation}

When the hidden-sector influence is stable and progressively screened, the leading memory kernel is expected to take an exponentially damped form,
\begin{equation}
K(t,t')
\propto
\exp[-\lambda(t-t')]\Theta(t-t').
\end{equation}

This provides the theoretical bridge between the foundational SBC postulate of finite observational access and the phenomenological SBC/F3 infrared response.
